% ------------------------------------------------------------
% Page 84
% ------------------------------------------------------------

\textbf{La tragédie de LA PARADE}

\vspace{0.35cm}

Le dimanche avant Pentecôte, jour de la fête du Rocher de la Vierge, le maquis Bir Hakeim
avait envahi Meyrueis et désarmé les gendarmes. Le jeune Noguès qui avait 17 ans suivit le
maquis. Je lui avais réglé son compte à l’usine la veille. René Fages (dit Fafa) à qui j’avais
fait la même chose partit également. Le malheureux Noguès une semaine après était pris et
fusillé au col de la Tourette. Fages plus heureux avait réussi à s’échapper (en compagnie de
Pierre Damiani dit Popeye.) Rolland et le pasteur Robert purent les récupérer, blessés,
dans la vallée de la Jonte (tous deux avaient réussi à se traîner jusqu’au Maynial car
Damiani avait été atteint par une rafale de mitraillette au bas du dos) et les planquer au
Moulin d’Ayres près de Meyrueis sur la route de Florac chez Mme. Loupiac.
Notre ami le Dr. Bouisset vint les soigner car Fages avait lui aussi deux balles dans la cuisse.
(En ce qui me concerne je n’ai gardé le souvenir que du Dr. Juliette Caban el venant, dès leur
arrivée (?) leur faire des piqûres antitétaniques tout en ne réussissant pas à extraire une ou
deux balles de leurs blessures. Ils étaient couchés dans la chambre de mes parents au 1er
Etage et chaque jour (?) je venais leur apporter de quoi manger en m’approchant
discrètement faisant semblant de pêcher, car le Moulin restait fermé. Ils restèrent quelques
jours, mais après le passage de la dernière colonne allemande, Rolland et le Pasteur Robert
vinrent les prendre avec la camionnette gazogène pour les cacher chez le Pasteur Chazel à
Vébron. *

\vspace{0.45cm}

Le 6 Juin 1944 une colonne allemande d’une trentaine de véhicules arrivait par la route de
Florac et s’arrêta à l’entrée de Meyrueis (mettant en batterie un petit canon au croisement de
la Fabrique) Sylvain Rey arriva en trombe dans mon bureau disant : « ça y est ! les
américains ont débarqué en Normandie » Sa femme, de l’hôtel, par téléphone, donna l’alerte
concernant l’arrivée des allemands. Aussitôt les irréguliers se retirèrent vers les bois de
Raffègues.
Puis une estafette venue certainement de Mende leur apporta un contre ordre et la colonne
repartit sans avoir exécuté sa mission qui, après ce qui a été dit ensuite devait consister à
arrêter les autorités municipales, à fouiller et à incendier les maisons pour rechercher « les
terroristes »échappés du maquis de la Parade ». Meyrueis l’a échappé belle. (Nous vîmes à
travers les fentes des volets, Damiani, Fafa et moi passer les camions sur la route au dessus
du moulin, s’arrêtant même pour une pose pipi. Quelle frayeur ! ! *

\vspace{0.55cm}

\textbf{Les départs au maquis.}

\vspace{0.35cm}

Peu après ces événements, un matin, en arrivant à l’usine j’appris le départ de Rolland pour le
maquis de l’Aigoual. S’étaient joints à lui son frère « Bonnet », le chef mécanicien ainsi que
4 ou 5 autres réfractaires de l’usine. nos amis le Dr. Bouisset et M. Griscelli, percepteur,
étaient également partis. Restaient pour faire tourner l’usine qui marchait au ralenti M.
Charles, moi même, et quelques ouvriers.

\vspace{0.45cm}

Le maquis vint nous prendre le dernier camion appelé « l’oiseau bleu », le seul véhicule qui
pouvait encore ravitailler Meyrueis. Grâce à une intervention de M. Charles qui put faire
toucher Rolland à l’Aigoual le camion nous fut rendu.

\vspace{0.45cm}

En tout et pour tout nous avons été obligés de livrer à la Kommandantur de Mende, un camion
de charbon de bois. Le chargement avait été tellement bien effectué que dans les virages de la
